\documentclass[11pt, reqno]{article}

\usepackage{amsmath, amsthm, amssymb}
\usepackage{enumitem}
\usepackage{tcolorbox}
\usepackage{hyperref}
\usepackage{tikz}
\usepackage{tikz-cd}
\usepackage{pgfplots}
\pgfplotsset{compat=1.18}
\usetikzlibrary{arrows.meta}
\usepackage{mathrsfs}
\usepackage{fancyhdr}
\usepackage[bottom=0.75in, top=1in, left=0.5in, right=0.5in]{geometry}
\usepackage{array}   % for \newcolumntype macro
\newcolumntype{L}{>{$}l<{$}}

\theoremstyle{plain}
\newtheorem*{theorem}{Theorem}
\newtheorem*{proposition}{Proposition}
\newtheorem{exercise}{Exercise}
\newtheorem*{lemma}{Lemma}
\newtheorem*{corollary}{Corollary}

\theoremstyle{definition}
\newtheorem*{definition}{Definition}
\newtheorem*{example}{Example}

\theoremstyle{remark}
\newtheorem*{remark}{Remark}

\renewcommand{\phi}{\varphi}
\renewcommand{\epsilon}{\varepsilon}
\renewcommand{\emptyset}{\varnothing}

\newcommand{\RR}{\mathbb{R}}
\newcommand{\ZZ}{\mathbb{Z}}
\newcommand{\NN}{\mathbb{N}}
\newcommand{\CC}{\mathbb{C}}
\newcommand{\QQ}{\mathbb{Q}}

\DeclareMathOperator{\ima}{\text{im}}

\begin{document}

\topmargin=-40pt
\rhead{Henry Woodburn}
\renewcommand{\headrulewidth}{1pt}
\renewcommand{\headsep}{20pt}
\thispagestyle{fancy}

{\Huge \bfseries \noindent Hitchin Harmonic Spinors Notes}

\section{Outline}
\begin{enumerate}
    \item \textbf{Preliminaries}
    \begin{enumerate}
        \item Clifford algebra
        \item spin groups and representation from last time
        \item lifting connection to spin
        \item spin c
    \end{enumerate}

    \item theorem statement and proof 
\end{enumerate}

Let $V$ be an even dimensional spin vector bundle. If we equip $V$ 
with a connection $\nabla$, we get a corresponding connection on the principal-$\text{SO}(2n)$
bundle of frames of $V$, denoted $P_{\text{SO}}$. This lifts to a connection on $P_{\text{Spin}}$, 
the principal-$\text{spin}(2n)$ bundle on $V$, by the covering map. 

One can show that if the connection $1$-form is given locally by $\omega_{ij}$, then relative
to this lifting, the lifted matrix is given by 
\[
    \frac{1}{4}\sum_{i,j}\omega_{ij}e_i e_j
\]

 

\end{document}